% !TEX program = lualatex
\documentclass[11pt]{memo_template}

\memonum{TCM-YYYY-XXXX}
\title{Buzz words without any significant meaning.  A really long
title may wraps onto another line}
\author{Your least appropriate childhood nickname}
\date{2461250.5 JD}
\memoto{Someone who might read your memo}
\memocc{Someone who probably won't read your memo
  \newline Someone who definitely won't read your memo
\newline Someone who will disagree with your memo before fully reading it}
\memosup{Nothing or everything}

\usepackage{pdfpages}

\begin{document}
\includepdf{coverpage.pdf}
\clearpage
\setcounter{page}{1}

\maketitle

\section{Introduction}
This is a simple article template with 1-inch margins on letter-size paper.

\blindtext

\section{Background}
\Blindtext

\section{Methods}
\blindtext

The best memos have some math,
\begin{equation}
  \left[-\frac{\hbar^2}{2\mu}\nabla^2 - \frac{e^2}{4\pi\varepsilon_0
  r}\right]\psi(r,\theta,\phi)=E\psi(r,\theta,\phi).
\end{equation}

Good memos might just have some code,
\begin{verbatim}
  % MATLAB sample
  x = linspace(0, 2*pi, 100);
  y = sin(x);
  plot(x, y)
  axis([0 2*pi -1.2 1.2])
\end{verbatim}

\section{Conclusion}
\blindtext

\includepdf{coverpage.pdf}
\end{document}
